% ============================================================================
% MUSTERLÖSUNG: Reibungskräfte (UE 5)
% Physik Klasse 7 - Gesamtschule MV
% ============================================================================
\documentclass[a4paper,11pt]{article}

\usepackage[left=1.5cm,right=1.5cm,top=1.5cm,bottom=1.5cm]{geometry}
\usepackage[utf8]{inputenc}
\usepackage[T1]{fontenc}
\usepackage[ngerman]{babel}
\usepackage{tikz}
\usepackage{fancyhdr}
\usepackage{tcolorbox}
\usepackage{amssymb,amsmath}
\usepackage{xcolor}
\usepackage{tabularx}
\usepackage{colortbl}
\usepackage{qrcode}
\usepackage[hidelinks]{hyperref}

\renewcommand{\familydefault}{\sfdefault}

% Seiteneinstellungen
\pagestyle{fancy}
\fancyhf{}
\renewcommand{\headrulewidth}{0.4pt}
\lhead{\textbf{Physik Klasse 7}}
\chead{Reibungskräfte -- \textcolor{red}{MUSTERLÖSUNG}}
\rhead{}
\setlength{\headheight}{14pt}
\cfoot{\thepage}

% Farben
\definecolor{physik}{HTML}{2c5aa0}
\definecolor{success}{HTML}{2a7a4b}
\definecolor{warning}{HTML}{b35c00}
\definecolor{loesung}{HTML}{c62828}

% Hilfsbefehle
\newcommand{\luecke}[1]{\underline{\hspace{#1}}}
\newcommand{\stern}[1]{\textcolor{warning}{\textbf{#1}}}
\newcommand{\lsg}[1]{\textcolor{loesung}{\textbf{#1}}}

\setlength{\parindent}{0pt}
\setlength{\parskip}{6pt}

\begin{document}

% ==================== KOPFBEREICH K0 ====================

\begin{tcolorbox}[colback=white, colframe=physik, boxrule=1.5pt, arc=0pt,
    title={\textbf{Reibungskräfte -- \textcolor{red}{MUSTERLÖSUNG}}}, fonttitle=\large, coltitle=black, colbacktitle=white]

\begin{tabular}{@{}ll@{\hspace{2cm}}ll@{}}
\textbf{Klasse:} & \lsg{7A/7B} & \textbf{Datum:} & \lsg{\_\_\_\_\_\_\_\_\_\_}
\end{tabular}

\end{tcolorbox}

\vspace{3mm}

% ==================== INFOKASTEN K1 ====================

\begin{tcolorbox}[colback=white, colframe=physik, boxrule=1pt, arc=0pt,
    title={\textbf{K1: Was ist Reibung?}}, fonttitle=\normalsize, coltitle=black, colbacktitle=white]
\textbf{Reibung} ist eine Kraft, die einer Bewegung entgegenwirkt. Sie entsteht, wenn zwei Oberflächen aneinander vorbeigleiten oder gleiten wollen.
\end{tcolorbox}

\vspace{3mm}

% ==================== BEOBACHTUNG B1 ====================

\textbf{B1: Beobachtung -- Klotz anschieben} \hfill (2 BE)

Um den Klotz in Bewegung zu setzen, muss die Kraft erst einen bestimmten \lsg{Wert / Schwellwert} überschreiten.

Sobald der Klotz gleitet, ist die benötigte Kraft \lsg{kleiner} (größer / kleiner).

\vspace{3mm}

% ==================== TABELLE T1 ====================

\textbf{T1: Messwerte -- Drei Reibungsarten} \hfill (3 BE)

\begin{center}
\begin{tabular}{|l|c|c|}
\hline
\rowcolor{white}
\textbf{Reibungsart} & \textbf{Kraftmesser-Wert} & \textbf{Symbol} \\
\hline
Haftreibung (Losbrechkraft) & \lsg{8} N & $F_{\text{Haft}}$ \\
\hline
Gleitreibung & \lsg{5} N & $F_{\text{Gleit}}$ \\
\hline
Rollreibung & \lsg{1} N & $F_{\text{Roll}}$ \\
\hline
\end{tabular}
\end{center}

\vspace{3mm}

% ==================== BEOBACHTUNG B2 ====================

\textbf{B2: Reihenfolge der Reibungskräfte} \hfill (2 BE)

Setze $>$ oder $<$ ein:

\begin{center}
\Large
$F_{\text{Haft}}$ \lsg{$>$} $F_{\text{Gleit}}$ \lsg{$>$} $F_{\text{Roll}}$
\end{center}

\vspace{3mm}

% ==================== MERKKASTEN K2 ====================

\begin{tcolorbox}[colback=white, colframe=physik, boxrule=1.5pt, arc=0pt,
    title={\textbf{K2: Die drei Reibungsarten}}, fonttitle=\normalsize, coltitle=black, colbacktitle=white]

\begin{tabular}{@{}p{3.5cm}p{10cm}@{}}
\textbf{Haftreibung:} & Kraft, um einen ruhenden Körper in Bewegung zu setzen. \\[2mm]
& Sie ist am \lsg{größten}. \\[4mm]
\textbf{Gleitreibung:} & Kraft während der \lsg{Bewegung} eines Körpers. \\[2mm]
& Sie ist \lsg{kleiner} als die Haftreibung. \\[4mm]
\textbf{Rollreibung:} & Kraft bei \lsg{rollenden} Körpern (z.B. Räder). \\[2mm]
& Sie ist am \lsg{kleinsten}. \\
\end{tabular}

\end{tcolorbox}

\vspace{3mm}

% ==================== BEOBACHTUNGEN B3, B4 ====================

\textbf{B3 + B4: Einflussfaktoren} \hfill (4 BE)

\begin{tabular}{@{}p{7.5cm}p{7.5cm}@{}}
\textbf{B3: Einfluss der Gewichtskraft} & \textbf{B4: Einfluss der Oberfläche} \\[2mm]
Mehr Gewichtskraft bedeutet: & Rauere Oberfläche bedeutet: \\[2mm]
\lsg{größere} Reibungskraft & \lsg{größere} Reibungskraft \\
\end{tabular}

\vspace{3mm}

% ==================== TABELLE T2 ====================

\textbf{T2: Erwünschte und unerwünschte Reibung} \hfill (3 BE)

\begin{center}
\begin{tabular}{|p{6.5cm}|p{6.5cm}|}
\hline
\rowcolor{white}
\textbf{Erwünscht (nützlich)} & \textbf{Unerwünscht (störend)} \\
\hline
1. \lsg{Bremsen funktionieren} & 1. \lsg{Motor verschleißt} \\[3mm]
2. \lsg{Schuhe rutschen nicht} & 2. \lsg{Energie geht verloren} \\[3mm]
3. \lsg{Schrauben halten} & 3. \lsg{Gelenke nutzen ab} \\
\hline
\end{tabular}
\end{center}

\vspace{5mm}

% ==================== AUFGABEN STERN ====================

\begin{tcolorbox}[colback=white, colframe=warning, boxrule=1pt, arc=0pt,
    title={\textbf{Aufgaben \stern{*}}}, fonttitle=\normalsize, coltitle=black, colbacktitle=white]

\textbf{A1:} Beim Bremsen eines Autos: Welche Reibungsart ist wichtig? \hfill (2 BE)

\lsg{Haftreibung} (zwischen Reifen und Straße; solange das Rad nicht blockiert)

\vspace{2mm}

\textbf{A2:} Nenne je ein Beispiel am Fahrrad: \hfill (2 BE)

Reibung erwünscht: \lsg{Bremsen / Reifen auf Straße}

Reibung unerwünscht: \lsg{Kette / Lager / Tretlager}

\end{tcolorbox}

\vspace{5mm}

% ==================== AUFGABEN STERN STERN ====================

\begin{tcolorbox}[colback=white, colframe=warning, boxrule=1pt, arc=0pt,
    title={\textbf{Aufgaben \stern{**}}}, fonttitle=\normalsize, coltitle=black, colbacktitle=white]

\textbf{A3:} Beschreibe in eigenen Worten: Wovon hängt die Reibungskraft ab? \hfill (3 BE)

\lsg{Die Reibungskraft hängt ab von:}
\begin{itemize}
    \item \lsg{der Gewichtskraft / Normalkraft (je schwerer, desto mehr Reibung)}
    \item \lsg{der Oberfläche / Rauheit (je rauer, desto mehr Reibung)}
\end{itemize}

\end{tcolorbox}

\vspace{3mm}

% ==================== FORMELKASTEN K3 (***) ====================

\begin{tcolorbox}[colback=white, colframe=physik, boxrule=1.5pt, arc=0pt,
    title={\textbf{K3: Formel für die Reibungskraft \stern{***}}}, fonttitle=\normalsize, coltitle=black, colbacktitle=white]

\begin{center}
\Large
\fbox{\rule{0pt}{1.5em}\hspace{1em}$F_R = $ \lsg{$\mu \cdot F_N$}\hspace{1em}}
\end{center}

\vspace{2mm}

\begin{tabular}{@{}lll@{}}
$F_R$ & = & Reibungskraft in \lsg{N} \\[2mm]
$\mu$ & = & Reibungszahl (griechisch: \glqq my\grqq), \lsg{ohne Einheit} \\[2mm]
$F_N$ & = & Normalkraft in N (auf waagerechter Fläche = \lsg{Gewichtskraft}) \\
\end{tabular}

\vspace{3mm}

\textbf{Reibungszahlen (Beispiele):}

\begin{center}
\begin{tabular}{|l|c|}
\hline
\rowcolor{white}
\textbf{Materialpaarung} & \textbf{$\mu$} \\
\hline
Gummi -- Asphalt (trocken) & 0,8 \\
\hline
Holz -- Holz & 0,5 \\
\hline
Stahl -- Eis & 0,03 \\
\hline
\end{tabular}
\end{center}

\end{tcolorbox}

\vspace{5mm}

% ==================== AUFGABEN STERN STERN STERN ====================

\begin{tcolorbox}[colback=white, colframe=warning, boxrule=1pt, arc=0pt,
    title={\textbf{Aufgaben \stern{***}}}, fonttitle=\normalsize, coltitle=black, colbacktitle=white]

\textbf{A4:} Lies die Reibungszahlen aus der Tabelle ab und vergleiche: \hfill (3 BE)

a) Welche Materialpaarung hat die größte Reibung? \lsg{Gummi -- Asphalt}

b) Welche hat die kleinste Reibung? \lsg{Stahl -- Eis}

c) Warum ist $\mu$ bei Stahl--Eis so klein? \lsg{Eis ist sehr glatt / geringe Rauheit}

\textbf{A5:} Erkläre: Warum verringert ein Schmierfilm (Öl oder Wasser) die Reibung? \hfill (3 BE)

\lsg{Der Schmierfilm trennt die beiden Oberflächen voneinander. Die Flächen berühren sich nicht mehr direkt, sondern gleiten auf dem Ölfilm. Dadurch wird die Reibung (Gleitreibung) verringert.}

\end{tcolorbox}

\vspace{3mm}

% ==================== PUNKTEKASTEN ====================

\begin{center}
\begin{tabular}{|c|c|c|}
\hline
\rowcolor{white}
\textbf{\stern{*} Basis} & \textbf{\stern{**} Standard} & \textbf{\stern{***} Erweitert} \\
\hline
\lsg{21} / 21 BE & \lsg{24} / 24 BE & \lsg{30} / 30 BE \\
\hline
\end{tabular}
\end{center}

\end{document}
